% 第19节：开放问题与未来方向
\section{开放问题与未来方向}
\label{sec:open_problems}

\subsection{数学挑战}

尽管本文档记录了进展，但仍有若干重大数学挑战：

\subsubsection{III型因子结构}

\begin{problem}
证明CTUFT中局域代数$\mathcal{A}(\mathcal{O})$是III$_1$型因子并确立split性质。
\end{problem}

\textit{状态}：基于一般量子场论结果，预期具有III$_1$型性质，但对于挠率场的严格证明需要证明模谱为$\mathbb{R}_+$。

\subsubsection{非微扰构造}

\begin{problem}
使用格点正则化或构造性场论方法非微扰地构造相互作用Wightman函数。
\end{problem}

\textit{状态}：当前结果是微扰的。格点表述存在，但连续极限需要进一步分析。

\subsubsection{无限维测度理论}

\begin{problem}
为几何量子化中的无限维相空间发展严格的测度理论。
\end{problem}

\textit{状态}：有限模式截断是良定义的，但极限需要仔细处理无限维Wiener测度。

\subsection{物理扩展}

\subsubsection{量子信息论}

\begin{problem}
在代数框架中使用复制品技巧计算纠缠熵，并联系黑洞熵。
\end{problem}

\subsubsection{宇宙学应用}

\begin{problem}
将严格框架扩展到弯曲时空背景，特别是暴胀几何。
\end{problem}

\subsection{跨学科联系}

\subsubsection{非交换几何}

\begin{problem}
确立CTUFT与Connes非交换几何计划的精确关系。
\end{problem}

\subsubsection{弦理论}

\begin{problem}
研究CTUFT是否作为弦理论的低能极限涌现，或代表不同的UV完备。
\end{problem}

\subsection{实验意义}

这些数学问题的解决具有直接的物理后果：
\begin{itemize}
    \item III型结构$\to$真空的纠缠性质
    \item 非微扰效应$\to$散射振幅的修正
    \item 弯曲时空扩展$\to$宇宙学预测
\end{itemize}
